% !TeX root = ../../main.tex
% =====================================================================
%  Anexo D — Código
%  Los fuentes van en la carpeta Codigo/ y se insertan con \lstinputlisting.
%  NUNCA pegues código a mano: se queda desactualizado en cuanto lo tocas.
% =====================================================================
\chapter{Código}
\label{ch:codigo}

Los programas del equipo se guardan en \texttt{Codigo/} y se insertan desde el
fichero fuente, nunca copiados a mano: así el libro y el código no se separan.

\begin{verbatim}
\lstinputlisting[language=Matlab,
                 caption={Cálculo de la eficacia del radiador (método NTU)},
                 label={lst:ntu}]{refrigeracion/EvaluaConfiguracion.m}

% Solo unas líneas concretas:
\lstinputlisting[language=Matlab, firstline=40, lastline=75,
                 caption={Bucle de barrido de configuraciones}]
                {refrigeracion/MAIN.m}
\end{verbatim}

\noindent
La ruta es relativa a \texttt{Codigo/} (está fijada con
\verb|\lstset{inputpath=Codigo/}| en el preámbulo).

\section{Refrigeración}
\label{sec:codigo_refrigeracion}
\pendiente{Modelo de radiador en MATLAB: MAIN, Config, EvaluaConfiguracion y test de regresión.}

\section{Dinámica vehicular}
\label{sec:codigo_dinamica}

\begin{description}[leftmargin=0pt,labelsep=0pt,style=nextline]
  \item[\texttt{dinamica/ajuste\_mf.py}] Ajuste de la Magic Formula a datos de
    ensayo de neumático: agrupa por carga, ajusta escalón a escalón, comprueba
    que el resultado tenga sentido físico y dibuja ajuste y residuos. Sin
    fichero de entrada genera datos sintéticos, de modo que el flujo se puede
    probar antes de tener datos reales. Se explica en el capítulo
    \ref{ch:modelado_neumatico} (códigos \ref{lst:mf_modelo} y
    \ref{lst:mf_ajuste}). Requiere \texttt{numpy}, \texttt{scipy} y
    \texttt{matplotlib}.
\end{description}

\pendiente{Faltan la simulación de vuelta y la generación del diagrama YMD.}

\section{Suspensión}
\label{sec:codigo_suspension}
\pendiente{Cálculo de cinemática, relaciones de instalación y casos de carga.}

\section{Utilidades}
\label{sec:codigo_utilidades}
\pendiente{Scripts de tratamiento de datos de ensayo.}

